\subsection*{Пример 1}

Исходное уравнение на всей действительной прямой:

\[a(x)y^{\prime\prime}(x)+b(x)y^\prime(x)+c(x)y(x)=g(x),\qquad x\in\mathbb R.\]

Все числа ниже сохраняют числовые параметры исходной задачи. Вспомогательные функции полностью определены последовательно; неопределённых коэффициентов или подбираемых параметров нет.

Для компактной записи аргумент x у вспомогательных функций опущен.

\[p_{1}=-0.9445099704819986.\]

\[q_{1}=-1.4320384163375453.\]

\[r_{1}=-1.0418788348584225.\]

\[h_{1}=e^{- 1.6583964568669665 \left(x + 0.7759412582493684\right)^{2}}.\]

\[d_{1}=\left(- 3.316792913733933 x - 2.5736364668352967\right) e^{- 1.6583964568669665 \left(x + 0.7759412582493684\right)^{2}}.\]

\[e_{1}=\left(\left(3.316792913733933 x + 2.5736364668352967\right)^{2} - 3.316792913733933\right) e^{- 1.6583964568669665 \left(x + 0.7759412582493684\right)^{2}}.\]

\[h_{2}=0.774948700390804.\]

\[d_{2}=0.\]

\[e_{2}=0.\]

\[h_{3}=1.4204653788060875 h_{1} + 1.4499935751532103 h_{2}.\]

\[d_{3}=1.4204653788060875 d_{1} + 1.4499935751532103 d_{2}.\]

\[e_{3}=1.4204653788060875 e_{1} + 1.4499935751532103 e_{2}.\]

\[a(x)=0.4948927992340211 \tanh{\left(p_{1} \right)} + 1.983925129927827,\quad b(x)=1.2736839759694616 \tanh{\left(q_{1} \right)},\quad c(x)=1.2336798048825266 \tanh{\left(r_{1} \right)} - 1.6551040450654098.\]

\[z=\frac{x-(0.12552944688615264)}{2.7787272448026363},\qquad D=1-z^2,\qquad s=0.4496376326792234.\]

\[W(x)=\begin{cases}\exp(1-D^{-1}),&|z|<1,\\0,&|z|\ge1,\end{cases}\qquad y(x)=sW(x)h_{3}.\]

\[\ell(x)=-\frac{2z}{2.7787272448026363D^2},\qquad \ell_1(x)=-\frac{2}{(2.7787272448026363)^2D^2}-\frac{8z^2}{(2.7787272448026363)^2D^3}.\]

\[g(x)=\begin{cases}s e^{1-D^{-1}}\left\{a(x)e_{3}+[2a(x)\ell(x)+b(x)]d_{3}+[a(x)(\ell(x)^2+\ell_1(x))+b(x)\ell(x)+c(x)]h_{3}\right\},&|z|<1,\\0,&|z|\ge1.\end{cases}\]

Здесь $d_{3}$ и $e_{3}$ заданы явными формулами выше и равны первой и второй производным $h_{3}$. Таким образом, в правой части нет невычисленных производных.

Границы носителя: -2.6531977979164836 и 2.904256691688789. Гладкое продолжение решения и правой части за носителем равно нулю; коэффициенты определены на всей прямой.

Резольвентная постановка: $L=a(x)\partial_x^2+b(x)\partial_x+c_R(x)$, где

\[\lambda=1.6551040450654098,\quad c_R(x)=c(x)+1.6551040450654098,\quad g_R(x)=-g(x),\qquad(\lambda-L)y=g_R.\]

\[\underline a=1.489032330693806,\quad\overline a=2.478817929161848,\quad |b|\le1.2736839759694616,\quad |c_R(x)|\le C=1.2336798048825266,\quad\lambda-C\approx0.42142424018288316>0.\]

Оценки $\underline a\le a(x)\le\overline a$, $|b(x)|$ и $C$ глобальны и могут быть неточными; они не названы точными экстремумами.

Ищется ограниченное решение на всей прямой. Отрезок [-1,1] - область сравнения сохранённых численных значений, а не граница определения коэффициентов.

Отдельная совместимая задача Коши на $[-1,1]$ имела бы дополнительные условия

\[y(-1)\approx0.8981507364381763,\qquad y^\prime(-1)\approx0.7336190907516743.\]

Эти значения следуют из указанного точного решения. Они не заменяют условие ограниченности в резольвентной задаче; произвольные другие начальные данные определяли бы другую задачу Коши.
